\subsection{不可分解内射模}

设 A 为环。关系「I 是不可分解内射 A-模的一个类」是集化关系（A, X, 第 21 页，推论 1）；我们将用 \(\mathcal{I}(A)\) 表示不可分解内射 A-模的诸类所成的集合。

命题 1.—— 设 A 为 Noether 环。对 A 的每个素理想 p，设 \(e_{\mathfrak{p}}: \mathrm{A}/\mathfrak{p} \to \mathrm{I}(\mathfrak{p})\) 为 \(\mathrm{A}/\mathfrak{p}\) 的一个内射包（A, X, § 1, 第 9 节）。

\begin{enumerate}
\def\labelenumi{\alph{enumi})}
\item
  A-模 I(p) 都是不可分解的。
\item
  设 I 为不可分解内射 A-模；集合 Ass(I) 只含一个元素。
\item
  映射 \(\mathfrak{p} \mapsto \operatorname{cl}(\operatorname{I}(\mathfrak{p}))\) 是从 \(\operatorname{Spec}(\operatorname{A})\) 到 \(\mathcal{I}(A)\) 上的双射。其逆双射把 \(\mathcal{I}(\operatorname{A})\) 的元素 I 对应到 \(\operatorname{Ass}(\operatorname{I})\) 的唯一元素。
\end{enumerate}

设 \(\mathfrak{p} \in \operatorname{Spec}(A)\)；我们来证明模 \(I(\mathfrak{p})\) 不可分解。由 A, X, 第 21 页，推论 2，只需证明：若 \(\mathfrak{a}\) 与 \(\mathfrak{b}\) 是 A 的包含 \(\mathfrak{p}\) 且不等于 \(\mathfrak{p}\) 的理想，则理想 \(\mathfrak{a} \cap \mathfrak{b}\) 不等于 \(\mathfrak{p}\)；或者，若 \(a\) 是 \(\mathfrak{a} - \mathfrak{p}\) 的元素而 \(b\) 是 \(\mathfrak{b} - \mathfrak{p}\) 的元素，则乘积 \(ab\) 属于 \((\mathfrak{a} \cap \mathfrak{b}) - \mathfrak{p}\)。

设 I 为不可分解内射 A-模，p、q 为 Ass(I) 的元素。于是 I 包含一个同构于 A/p 的子模 M 和一个同构于 A/q 的子模 N。我们有 M ∩ N ≠ 0（A, X, 第 21 页，命题 14）；对 M ∩ N 的每个非零元素 x，有 p = Ann x = q。由于 Ass(I) 非空（IV, § 1, 第 1 节，命题 2 的推论 1），它只含一个元素 p(I)。

这样我们定义了两个映射：\(\mathfrak{p} \mapsto \operatorname{cl}(I(\mathfrak{p}))\)（从 \(\operatorname{Spec}(A)\) 到 \(\mathcal{I}(A)\)）与 \(I \mapsto \mathfrak{p}(I)\)（从 \(\mathcal{I}(A)\) 到 \(\operatorname{Spec}(A)\)）；我们来证明这两个映射互为逆双射。设 \(\mathfrak{p} \in \operatorname{Spec}(A)\)；则 \(\mathfrak{p}\) 属于 \(\operatorname{Ass}(I(\mathfrak{p}))\)，因而是 \(\operatorname{Ass}(I(\mathfrak{p}))\) 的唯一元素。设 \(I\) 为不可分解内射 A-模，而 \(\mathfrak{p}\) 为 \(\operatorname{Ass}(I)\) 的唯一元素；则 \(I\) 是 \(A/\mathfrak{p}\) 的一个内射包（A, X, 第 21 页，命题 14）。命题证毕。

注 1.——设 I 为不可分解内射 A-模，p 为 Ass(I) 的唯一元素；由命题 1，I 包含一个同构于 A/p 的子模，且 I 是它的内射包。一般而言这样的子模并不唯一，取 A = Z、p = 0、I = Q 即见。

对 \(\mathfrak{p}\)（\(\mathrm{A}\) 中的素理想）的每一个，如上选取 \((\mathrm{I}(\mathfrak{p}), e_{\mathfrak{p}})\) 作为 \(\Lambda/\mathfrak{p}\) 的内射包。由 \(\mathrm{A}\), X, 第 22 页，定理 3，我们有：

定理 1.—— 设 A 为 Noether 环。对每个内射 A-模 I，存在唯一的基数族 \((a_{\mathfrak{p}})_{\mathfrak{p}\in\operatorname{Spec}(\mathrm{A})}\)，使 I 同构于 \(\bigoplus_{\mathfrak{p}}\mathrm{I}(\mathfrak{p})^{(a_{\mathfrak{p}})}\)。

由 IV, § 1, 第 1 节，命题 3 的推论 1，集合 Ass(I) 是族 \((a_{\mathfrak{p}})\) 的支撑集。

注 2. 设 \(A\) 为 Noether 环，M 为 A-模，\(e: \mathrm{M} \to \mathrm{I}\) 为 M 的一个内射包。集合 Ass(M) 等于 Ass(I)：事实上，包含关系 Ass(M) ⊂ Ass(I) 是显然的。另一方面，若 \(\mathfrak{p}\) 是 Ass(I) 的元素，则 I 包含一个同构于 \(A/\mathfrak{p}\) 的子模 N；由于 \(A\)-模 \(e^{-1}(\mathrm{N})\) 非零，我们有 \(\operatorname{Ass}(e^{-1}(\mathrm{N})) = \{\mathfrak{p}\}\)。这证明 \(\mathfrak{p}\) 与 \(e^{-1}(\mathrm{N})\) 相伴，从而与 M 相伴，断言得证。

采用定理中的记号，并进而假设 Ass(M) 有限。对每个 \(\mathfrak{q} \in \operatorname{Ass}(M)\)，用 \(Q_{\mathfrak{q}}\) 表示 \(\bigoplus_{\mathfrak{p} \in \operatorname{Ass}(M) - \{\mathfrak{q}\}} I(\mathfrak{p})^{(a_{\mathfrak{p}})}\) 与 M 的交。则 \((Q_{\mathfrak{q}})_{\mathfrak{q} \in \operatorname{Ass}(M)}\) 是 M 中零子模的一个既约准素分解（IV, § 2, 第 3 节，定义 3）。事实上，我们有 \(\cap Q_{\mathfrak{q}} = 0\)；由于 \(M/Q_{\mathfrak{q}}\) 等同于 \(I(\mathfrak{q})^{(a_{\mathfrak{q}})}\) 的一个非零子模，有 \(\operatorname{Ass}(M/Q_{\mathfrak{q}}) = \{\mathfrak{q}\}\)，于是只需应用同处的命题 4。

例.—— 设 A 为主理想环，K 为其分式域。内射 A-模就是可除 A-模（A, X, 第 17 页，推论 2）。A-模 K 是 A 的一个内射包（A, X, 第 20 页，例 1）。设 \(p\) 为 A 的极端元，\(\mathfrak{p}\) 为（极大）理想 \(Ap\)。用 \(e: A/\mathfrak{p} \longrightarrow K/A_{\mathfrak{p}}\) 表示把元素 \(a \in A\) 的类映到 \(a/p\) 的类的同态。我们来证明 \((K/A_{\mathfrak{p}}, e)\) 是 \(A/\mathfrak{p}\) 的一个内射包。同态 \(e\) 是单射。A-模 \(K/A_{\mathfrak{p}}\) 是可除模的一个商，故可除。设 \(x\) 为 \(K/A_{\mathfrak{p}}\) 的一个非零元素；它是元素 \(a/p^n\)（属于 K）的类，其中 \(a \in A - \{0\}\) 且 \(n \geqslant 1\)（A, VII, 第 10 页，定理 2）。于是我们有 \(p^{n-1}x = e(a)\)，从而 \(e(Ax) \neq 0\)，这就证明了我们的断言。

于是由定理 1，每个可除 A-模都是若干同构于 K 或 \(K/A_{p}\)（p 为 A 的某个极大（主）理想）的 A-模的直和。

